\documentclass[12pt]{article}%
\usepackage{amsfonts}
\usepackage{fancyhdr}
\usepackage{comment}
\usepackage[a4paper, top=2.5cm, bottom=2.5cm, left=2.2cm, right=2.2cm]%
{geometry}
\usepackage{times}
\usepackage{amsmath}
\usepackage{changepage}
\usepackage{amssymb}
\usepackage{graphicx}
\setcounter{MaxMatrixCols}{30}
\newtheorem{theorem}{Theorem}
\newtheorem{acknowledgement}[theorem]{Acknowledgement}
\newtheorem{algorithm}[theorem]{Algorithm}
\newtheorem{axiom}{Axiom}
\newtheorem{case}[theorem]{Case}
\newtheorem{claim}[theorem]{Claim}
\newtheorem{conclusion}[theorem]{Conclusion}
\newtheorem{condition}[theorem]{Condition}
\newtheorem{conjecture}[theorem]{Conjecture}
\newtheorem{corollary}[theorem]{Corollary}
\newtheorem{criterion}[theorem]{Criterion}
\newtheorem{definition}[theorem]{Definition}
\newtheorem{example}[theorem]{Example}
\newtheorem{exercise}[theorem]{Exercise}
\newtheorem{lemma}[theorem]{Lemma}
\newtheorem{notation}[theorem]{Notation}
\newtheorem{problem}[theorem]{Problem}
\newtheorem{proposition}[theorem]{Proposition}
\newtheorem{remark}[theorem]{Remark}
\newtheorem{solution}[theorem]{Solution}
\newtheorem{summary}[theorem]{Summary}
\newenvironment{proof}[1][Proof]{\textbf{#1.} }{\ \rule{0.5em}{0.5em}}

\newcommand{\Q}{\mathbb{Q}}
\newcommand{\R}{\mathbb{R}}
\newcommand{\C}{\mathbb{C}}
\newcommand{\Z}{\mathbb{Z}}

%%%%
% ME %
%%%%

\usepackage{listings}
\usepackage{color}
\usepackage{xcolor}
\usepackage{mdframed}

\definecolor{dkgreen}{rgb}{0,0.6,0}
\definecolor{gray}{rgb}{0.5,0.5,0.5}
\definecolor{mauve}{rgb}{0.58,0,0.82}

\lstset{%frame=tb,
language=Matlab,
aboveskip=3mm,
belowskip=3mm,
showstringspaces=false,
columns=flexible,
basicstyle={\small\ttfamily},
numbers=none,
numberstyle=\tiny\color{gray},
keywordstyle=\color{blue},
commentstyle=\color{dkgreen},
stringstyle=\color{mauve},
breaklines=true,
breakatwhitespace=true,
tabsize=3
}

\usepackage{outlines}
\usepackage{enumitem}
%\setenumerate[1]{label=\Roman*.}
%\setenumerate[2]{label=\Alph*.}
%\setenumerate[3]{label=\roman*.}
%\setenumerate[4]{label=\alph*.}

\usepackage{titlesec} 

\titleformat{\subsection}[runin]
  {\normalfont\large\bfseries}{\thesubsection}{1em}{}
\titleformat{\subsubsection}[runin]
  {\normalfont\normalsize\bfseries}{\thesubsubsection}{1em}{}

% for enumerate spacing
\newenvironment{tight_enum}{
\begin{enumerate}[label=\alph]
\setlength{\itemsep}{0pt}
\setlength{\parskip}{0pt}
}{\end{enumerate}}

\usepackage{exercise}

\newcommand{\code}[1]{\texttt{#1}}
\newcommand{\shp}{$>>>$ }
\newcommand{\tab}{\hspace*{22pt}}

\setlength\parindent{0pt}
\renewcommand{\theenumi}{\alph{enumi}}

\usepackage{multicol}
\setlength{\columnsep}{0.3cm}

\begin{document}

\title{Test 1}

\author{EEC 3150}
\date{02-14-23}

\maketitle

\begin{center}
\fbox{\fbox{\parbox{5.5in}{\centering
For each exercise, write a separate Python script. Submit all scripts on Blackboard under the appropriate assignment. Name your script files with the following format: \\
\vspace{10pt}
Test$<$test number$>$\_Ex$<$exercise number$>$\_$<$FirstInitial$><$LastName$>$.py \\
\vspace{10pt}
Example: Test01\_Ex01\_GWest.py
}}}
\end{center}


\pagebreak


\begin{multicols*}{2}

%%% NEW EXERCISE %%%
\begin{Exercise}[origin={25 points}]

\begin{enumerate}
\item Write a function \code{Force(q1,q2,r)} which will return the value of the force between two charged particles with charges $q_1,q_2$ and separation distance $r$. This can be calculated using Coulomb's Law
\begin{equation*}
F = k \frac{q_1 q_2}{r^2}
\end{equation*}
where $k = 9 \cdot 10^9$ (the units are ugly, so I omit them).
\item Test your function by using it to calculate the force between a proton and electron in a hydrogen atom: $q_{electon} = -1.6 \cdot 10^{-19} C, q_{proton} = 1.6 \cdot 10^{-19} C, r = 5.3 \cdot 10^{-11} m$.
\end{enumerate}


\end{Exercise}

%%% NEW EXERCISE %%%
\begin{Exercise}[origin={25 points}]

\begin{enumerate}
\item Write a function \code{HMS(n)} which will take an integer $n$ number of seconds and return a string that formats it as a number of hours, minutes, and seconds separated by spaces with units after the values. Here are some examples to illustrate how it should work: \\
\code{\shp HMS(30) \\
0h 0m 30s \\
\shp HMS(360) \\
0h 6m 0s \\
\shp HMS(3600) \\
1h 0m 0s \\
\shp HMS(4000) \\
1h 6m 40s
}
\item Demonstrate your function by calling it.
\end{enumerate}

\end{Exercise}

\vfill\null
\columnbreak

%%% NEW EXERCISE %%%
\begin{Exercise}[origin={25 points}]

\begin{enumerate}
\item Write a function \code{Dot(x,y)} which will return the dot product of two $n$-dimensional vectors $x,y$. Reminder:
\begin{equation*}
\vec x \cdot \vec y = \sum_{i=1}^n x_i y_i
\end{equation*}
\item Demonstrate your function by creating two lists and calling the function with them as input.
\end{enumerate}

\end{Exercise}

%%% NEW EXERCISE %%%
\begin{Exercise}[origin={25 points}]

\begin{enumerate}
\item Write a function \code{Mode(x)} which will return the mode (most frequently occurring element) of the list. Here is an example: \\
\code{\shp x = [1,1,2,2,2,3,3,3,3,3] \\
\shp Mode(x) \\
3
}
\item Demonstrate your function by creating a list calling the function with it.
\end{enumerate}

\end{Exercise}

\end{multicols*}




\end{document}






